\documentclass[11pt]{article}
\usepackage[margin=1in]{geometry}
\usepackage[T1]{fontenc}
\usepackage[utf8]{inputenc}
\usepackage{lmodern,microtype,amsmath,amssymb,mathtools}
\usepackage{enumitem,xcolor,booktabs}
\usepackage[colorlinks=true,linkcolor=black,citecolor=teal!60!black,urlcolor=teal!60!black]{hyperref}
\setlist{nosep,leftmargin=1.5em}
\newcommand{\SpanP}{\mathsf{SpanP}}
\newcommand{\UP}{\mathsf{UP}}
\title{Counting Fixed-Knot Completions in the Standard-Eleven Mosaic Model}
\author{Anonymous Authors}
\date{Manuscript prepared August 2026}
\begin{document}
\maketitle
\begin{abstract}
We study completion counting for the fixed eleven-tile
Lomonaco--Kauffman knot-mosaic alphabet.  An input is only a finite partial
assignment of ordinary microscopic tiles; every hole admits all eleven
tiles.  We define a polynomially recognizable crossing-blocked exact-core
class on which every legal one-component completion is automatically a
planar unknot.  By a solution-bijective reduction from maximum-degree-three
grid Hamiltonicity with marked forced edges, we prove that counting these
raw completions is \texttt{\#P}-complete under parsimonious reductions.  The proof is
microscopic: optional corridors are Boolean, vertex cells select a unique
degree-two mask, and completed mosaic components are exactly the cycles of a
spanning two-factor.

For every fixed classical knot \texttt{J}, we then replace one forced corridor by a
frozen two-terminal long-knot diagram of \texttt{J}.  Every legal completion becomes
the split union of \texttt{J} with planar unknots, without changing a single free
choice.  Counting completions of knot type \texttt{J} is therefore parsimoniously
\texttt{\#P}-complete in the raw unit-tile model.  On the same syntactic family,
Jones-polynomial equality with \texttt{V\_J}, determinant equality with \texttt{det(J)}, and
positivity of the Jones value at \texttt{-1} select exactly the one-component
completions and are likewise parsimoniously complete.  These are counting
results over exponentially many completions; evaluating \texttt{V\_L(-1)} for one
diagram remains polynomial-time.

We also identify the exact unrestricted boundary.  Exact-core raw unknot
counting is \texttt{\#P}-hard and belongs to \texttt{SpanP}; it belongs to \texttt{\#P} if and only if
mosaic unknot recognition belongs to \texttt{UP}.  Counting distinct blank-padded
rectangles is zero or infinite rather than a \texttt{\#P} function.  Finally, the
guarded functions are fixed-parameter tractable by hole count and board width,
while an FPRAS would imply \texttt{NP=RP}.
\end{abstract}
\section{Introduction}

Knot mosaics encode link diagrams by the eleven oriented local tiles
introduced by Lomonaco and Kauffman.  Existing work establishes equivalence
with tame-knot theory after stabilization, develops mosaic and tile numbers,
enumerates small mosaics, and implements SAT/ALLSAT and modern SageMath tools.
The question here is different: how hard is it to count completions of a
partially frozen microscopic mosaic when the completed diagram must have one
specified topology?

There are two traps.  First, a role-labelled macrotile reduction does not
automatically prove a theorem about arbitrary unit-cell choices.  Second,
unknot recognition being in \texttt{NP} does not automatically put the number of
distinct unknot completions in \texttt{\#P}: one diagram may have many accepting
topological certificates.  Our reduction avoids both issues.  Every hole is
a genuine unrestricted standard-eleven cell, while a local blocked-side guard
prevents all four-port tiles.  Thus topology collapses to a deterministic
component-count test on every input in the promised syntactic class.

The hardness source is the ASP-completeness theorem of the MIT Hardness Group
et al. for undirected maximum-degree-three grid Hamiltonicity.  ASP here is
the solution-bijection notion, not answer-set programming.  The constructive
proof supplies marked edges forced in every output solution.  Retaining those
marks gives exactly the local forcing needed at a mosaic vertex while
preserving the solution bijection.

The fixed-target extension is unusually clean.  A fixed long-knot tangle is
frozen into one already-forced corridor.  Replacing it by a boundary-parallel
arc preserves every component and every free tile choice.  With the tag
present, precisely one planar component is tied into \texttt{J}; all others remain
split unknots.  The result therefore covers every fixed classical knot and
also gives genuine Jones/determinant selectors without invoking the open
question whether the Jones polynomial detects the unknot.

\subsection{Contributions}

\begin{enumerate}
\item A parsimonious \texttt{\#P}-completeness theorem for raw standard-eleven
   one-component unknot completion on a polynomially recognizable domain.
\item A literal microscopic bijection with forced-edge spanning two-factors,
   including polynomial area, malformed inputs, and multiplicity one.
\item For every fixed knot \texttt{J}, a raw tagged theorem for knot type, Jones
   equality, determinant equality, and determinant positivity.
\item The exact \texttt{SpanP/UP} boundary for unrestricted exact-core unknot counting.
\item A proof that literal all-rectangle counting is zero or infinite.
\item Exact fixed-parameter algorithms and an FPRAS barrier.
\end{enumerate}

\section{Related work}

Lomonaco--Kauffman introduced the fixed eleven-tile model
(\href{https://doi.org/10.1007/s11128-008-0076-7}{2008}); Kuriya--Shehab proved its
stable equivalence with tame knots
(\href{https://doi.org/10.1142/S0218216514500035}{2014}).  Mosaic-number and
enumeration foundations include Lee--Hong--Lee--Oh
(\href{https://doi.org/10.1142/S0218216514500692}{2014}) and
Oh--Hong--Lee--Lee
(\href{https://doi.org/10.1007/s11128-014-0895-7}{2015}).  Miller's SAT and ALLSAT
work (\href{https://doi.org/10.1609/aaai.v36i11.21645}{2022},
\href{https://doi.org/10.1609/aaai.v37i13.27002}{2023}) is the closest direct
computational predecessor, but does not prove fixed-knot completion hardness.
Heap et al. provide a current low-crossing tabulation
(\href{https://doi.org/10.2140/involve.2025.18.91}{2025}); Deng et al. provide the
2026 KnotMosaics SageMath package (\href{https://arxiv.org/abs/2605.14189}{arXiv:2605.14189}).

Valiant introduced \texttt{\#P}
(\href{https://doi.org/10.1137/0208032}{1979}).  Köbler--Schöning--Torán introduced
span functions, counting the distinct outputs of an NP transducer
(\href{https://doi.org/10.1007/BF00276023}{1989}).  That is the appropriate upper
bound when an object is guessed together with a potentially nonunique
certificate.  Unknot recognition is in \texttt{NP} by Hass--Lagarias--Pippenger
(\href{https://doi.org/10.1145/301970.301971}{1999}); Lackenby gives efficient
knottedness certificates (\href{https://doi.org/10.1016/j.aim.2021.107796}{2021}).
Neither result supplies certificate uniqueness.

The source theorem is FUN 2024, Theorem 18
(\href{https://doi.org/10.4230/LIPIcs.FUN.2024.23}{DOI 10.4230/LIPIcs.FUN.2024.23});
the May 2026 full revision is \href{https://arxiv.org/abs/2405.08377}{arXiv:2405.08377}.
The latter updates authorship and puzzle applications; the proceedings result
is the peer-reviewed hardness source.

For the Jones consequences, Kauffman's state model supplies normalization
(\href{https://doi.org/10.1016/0040-9383(87}{1987}90009-7)), and Boninger gives a
recent direct Jones/Goeritz treatment of
\texttt{det(K)=|V\_K(-1)|}
(\href{https://doi.org/10.1112/blms.12753}{2023}).  Our result does not contradict
the easy evaluation at \texttt{-1}: it counts completions satisfying that predicate.

No predecessor satisfying all of the following axes was located in the
documented primary-source search through 14 August 2026: fixed standard-eleven
library, raw partial-matrix input, exact one-component fixed-knot acceptance,
and a parsimonious counting reduction.  This is a corpus-scoped search result,
not an absolute priority claim.

\section{Model and finite counting convention}

Let \texttt{T=\{0,...,10\}} be the oriented standard-eleven alphabet.  Tile \texttt{0} is
blank; tiles \texttt{1,...,6} have one noncrossing arc; tiles \texttt{7,8} have two disjoint
arcs; and \texttt{9,10} are the two crossings.  Rotations are distinct tile IDs.

A seed is a canonical finite partial map

\[
  s:\mathbb N^2\rightharpoonup T
\]

with nonnegative, duplicate-free coordinates.  A nonempty seed recovers

\[
 B_s=[0,h)\times[0,w),\quad
 h=1+\max\{r:(r,c)\in\operatorname{dom}s\},\quad
 w=1+\max\{c:(r,c)\in\operatorname{dom}s\}.
\]

A witness is one row-major tile value for every hole in \texttt{B\_s}.  Adjacent ports
must match and no occupied port may leave the board.  Strand components are
cycles of the resulting degree-two graph, keeping the two strands at a
crossing distinct.

A seed is \textbf{crossing-blocked} when:

\begin{enumerate}
\item it has no frozen crossing;
\item every hole has a side facing the exterior or a frozen blank neighbor;
\item \texttt{hw <= max(1,|dom(s)|)\^\{\}2};
\item every boundary cell is frozen blank; and
\item at least one interior cell is frozen nonblank.
\end{enumerate}

Malformed and non-guarded inputs have value zero.  This makes the counting
function total and the board polynomial in its encoding.

\subsection{Lemma — guarded normal form}

Every suitably connected filling of a crossing-blocked seed is crossing-free.
It is an unknot exactly when it has one strand component.

\textbf{Proof.}  Every four-port tile occupies the blocked side, so tiles \texttt{7--10}
are impossible at a hole; frozen crossings are forbidden.  A closed connected
one-manifold embedded in the plane is a simple closed curve. \hfill$\square$

\section{The parsimonious raw reduction}

Define \texttt{\#FGridHam\_3(G,F)} to count Hamiltonian edge sets of an undirected
maximum-degree-three grid graph \texttt{G} that contain the explicitly supplied set
\texttt{F}, where \texttt{F} is incident to every vertex.  Invalid instances have value
zero.  Retaining the marked black/tentacle edges in the constructive proof of
FUN 2024, Theorem 18, and marking an arbitrary incident edge at any remaining
degree-two vertex yields this total marked source without changing any
solution.  It is parsimoniously \texttt{\#P}-complete.

Translate a connected source graph to nonnegative coordinates and scale the
grid by four.  Freeze every internal cell of an edge in \texttt{F} to its unique path
tile.  Leave every source vertex and every internal cell of an optional edge
free.  Freeze all remaining cells blank and add a one-cell blank moat.  A
connected \texttt{n}-vertex unit-grid graph has coordinate span at most \texttt{n-1}, so the
compiled dense board has polynomial area even when source coordinates are
binary.

\subsection{Lemma — optional corridors}

An optional corridor is either entirely blank or entirely its unique path
tiles.

\textbf{Proof.}  At an internal corridor cell both transverse neighbors are frozen
blank.  Exactly two standard tiles are compatible: blank and the unique tile
joining predecessor to successor.  Port matching propagates that Boolean
choice along the whole corridor. \hfill$\square$

\subsection{Lemma — vertex selection}

Every suitable filling selects exactly two source edges incident to each
vertex and contains every edge in \texttt{F}.

\textbf{Proof.}  A forced corridor occupies at least one side.  Maximum degree
three leaves a blank-blocked side, excluding four-port tiles.  Every standard
tile has even occupied degree, so the free vertex cell has degree exactly two;
its two-side mask determines a unique tile. \hfill$\square$

\subsection{Lemma — microscopic bijection}

Suitable fillings are in bijection with spanning two-factors containing \texttt{F},
and mosaic components are the cycles of the two-factor.

\textbf{Proof.}  Lemmas 4.1--4.2 map a filling to a consistent selected edge set
of degree two.  Conversely, a two-factor uniquely sets every corridor to
blank or its path tile and uniquely sets each vertex tile.  The two maps are
inverse, and strand tracing follows graph cycles. \hfill$\square$

\subsection{Theorem — raw parsimonious completeness}

Let \texttt{\#CB-RKW11\_U(s)} count one-component unknot fillings of a crossing-blocked
exact core.  Then

\[
 \#\mathrm{CB\text{-}RKW11}_U(s_{G,F})
   =\#\mathrm{FGridHam}_3(G,F),
\]

and \texttt{\#CB-RKW11\_U} is \texttt{\#P}-complete under parsimonious reductions.

\textbf{Proof.}  By Lemma 4.3, one component is exactly one two-factor cycle, hence
a Hamiltonian cycle.  Lemma 3.1 makes it an unknot.  A \texttt{\#P} machine guesses
one four-bit code per hole, rejects codes \texttt{11--15}, and checks ports and
components deterministically.  Each matrix has exactly one accepted code. \hfill$\square$

\section{Every fixed target knot}

Fix an unoriented tame classical knot \texttt{J}.  Choose a diagram, remove a
noncrossing subarc adjacent to the unbounded face, rectilinearize the remaining
one-string tangle on the standard-eleven grid, and add a blank collar.  Denote
the resulting two-terminal matrix by \texttt{T\_J}.  Let \texttt{T\_0} be a boundary-parallel
arc on the same rectangle with the same two gates.  Both are constants because
\texttt{J} is fixed.

A seed is \textbf{\texttt{J}-guarded} when it has one fully frozen translated copy of
\texttt{T\_J}, every frozen crossing lies in that copy, every hole has a frozen blank
side, the area is quadratically bounded, and the outer boundary is a frozen
blank moat.  This is recognized by scanning constant-size translated
patterns; no knot-equivalence test occurs in the guard.

Replace the tag by \texttt{T\_0} without changing the holes.  The resulting completed
matrix is planar and has the same components.  Therefore every suitable
\texttt{J}-guarded completion with \texttt{k>=1} components has the complete topology

\[
 L(P)\cong J\sqcup U^{\sqcup(k-1)}.                 \tag{5.1}
\]

Indeed, the tag ties \texttt{J} into the unique planar component passing through its
two-point ball, while the other planar components remain split unknots.

To compile hardness, choose the first marked forced edge, rotate the entire
grid orientation-preservingly if needed, enlarge the constant grid scale, and
replace that forced corridor by \texttt{T\_J} and two frozen leads.  All optional
cells and vertex cells are exactly those of the untagged reduction.  Tag
erasure gives a component-preserving bijection with the original fillings.

\subsection{Theorem — raw fixed-target completeness}

For every fixed classical knot \texttt{J}, counting \texttt{J}-guarded raw standard-eleven
fillings of knot type \texttt{J} is \texttt{\#P}-complete under parsimonious reductions.
Hardness holds with all holes initially permitting all eleven tiles and with
only the constant number of crossings in the frozen tag.

\textbf{Proof.}  The source-to-filling bijection survives tag insertion.  By (5.1),
a filling has knot type \texttt{J} exactly when its two-factor has one cycle.  For
membership, guess each hole once and accept by deterministic local checks and
component count; (5.1) supplies the semantic equivalence. \hfill$\square$

\subsection{Corollary — genuine Jones and determinant selectors}

Let \texttt{V\_U(t)=1} and \texttt{delta(t)=-(t\^\{\}\{1/2\}+t\^\{\}\{-1/2\})}.  On every suitable
\texttt{J}-guarded completion,

\[
 V_{L(P)}(t)=\delta(t)^{k-1}V_J(t).                 \tag{5.2}
\]

Consequently the following four predicates are equivalent on this family:

\begin{enumerate}
\item \texttt{k=1};
\item \texttt{L(P)} has knot type \texttt{J};
\item \texttt{V\_\{L(P)\}(t)=V\_J(t)};
\item \texttt{|V\_\{L(P)\}(-1)|=det(J)}; equivalently \texttt{|V\_\{L(P)\}(-1)|>0}.
\end{enumerate}

If \texttt{k>1}, (5.2) differs from \texttt{V\_J} in the Laurent integral domain, and
\texttt{delta(-1)=0}; if \texttt{k=1}, equality is immediate.  Counting completions
satisfying any predicate is therefore parsimoniously \texttt{\#P}-complete.  A
Goeritz determinant computes the fourth predicate for one diagram in
polynomial time.  The theorem concerns counting all completing matrices.

\section{Unrestricted counting and padding}

Let \texttt{F\_raw(s)} count one-component unknot fillings of a polynomially bounded
exact core without the crossing-blocked syntax.  It is \texttt{\#P}-hard because the
guarded instances are ordinary raw instances.  It belongs unconditionally to
\texttt{SpanP}: an NP transducer guesses a matrix and an unknot certificate, verifies
them, and outputs only the matrix, so different certificates for the same
matrix are collapsed.

Moreover,

\[
 F_{raw}\in\#P
 \quad\Longleftrightarrow\quad
 \mathrm{UNKNOT}_{mosaic}\in UP.                  \tag{6.1}
\]

If \texttt{F\_raw} has a \texttt{\#P} machine, freeze every cell of a supplied closed mosaic;
its count is zero or one, giving an unambiguous unknot machine.  Conversely,
combine a canonical filling guess with an unambiguous unknot verifier.

If each larger rectangle is a distinct witness, any accepting core extends
by arbitrarily many blank rows and columns.  The literal all-size count is
therefore zero or infinite.  The exact-core convention is equivalent to a
quotient by exterior blank padding on guarded instances: the moat prevents a
strand crossing, and any nonblank exterior produces an additional component.

\section{Algorithms and approximation boundary}

With \texttt{h} holes, exact enumeration takes \texttt{O(11\^\{\}h poly(|B\_s|))} time.  After
erasing the fixed tag, a width-\texttt{w} scan stores the occupied frontier ports,
their noncrossing connectivity, and whether a component has closed.  There
are Catalan-many states, giving \texttt{|B\_s| 2\^\{\}\{O(w)\}} arithmetic operations.
Hardness already holds at crossing count zero for \texttt{J=U}, at the constant tag
crossing count for fixed \texttt{J}, and at target component count one.

A parsimonious reduction transfers any FPRAS to \texttt{\#SAT}.  SAT self-reduction,
amplification, and final witness verification then give an RP algorithm for
SAT.  Thus an FPRAS for any guarded fixed-target count implies \texttt{NP=RP}; a
deterministic multiplicative scheme that treats zero exactly implies \texttt{P=NP}.
No conclusion follows for a coarse additive error scaled by all \texttt{11\^\{\}h}
assignments, and no \texttt{BPP/BQP} collapse follows.

\section{Limitations}

The theorem is total on a natural, polynomially recognizable guarded family;
it does not establish \texttt{\#P} membership for unrestricted topological
completion.  It counts an exact recovered core or blank-padding equivalence
class, not every padded rectangle.  It does not require or imply a macrotile
parser.  Targets are unoriented classical knot types; framed, oriented,
satellite, concordance, and multi-component target classifications are not
claimed.  Jones equality is used only on the proved split-union family, not as
a general polynomial-time predicate.

\section{Reproducibility}

The raw baseline regressions exhaust all local source types and several
finite graphs.  The tagged regression replaces one forced edge of the marked
\texttt{2x4} ladder with a literal protected trefoil tangle, recovers the two expected
two-factors with component histogram \texttt{\{1:1,2:1\}}, verifies that tag erasure
preserves both component counts, and checks determinant three on the tagged
one-cycle example.  Corruption tests alter the tag, gate, moat, and source
marking and must be rejected.

The exact commands, environment, hashes, and evidence limits are recorded in
\texttt{08\_REPRODUCIBILITY.md}.  Finite tests guard implementation details; Lemmas
3.1 and 4.1--4.3 and the tag-erasure argument prove the unbounded statements.

\section{References}

The complete bibliography, including foundational work and the 2021--2026
search record, is in \texttt{07\_REFERENCES.bib}; the theorem-by-theorem inclusion and
exclusion matrix is \texttt{06\_LITERATURE\_MATRIX.csv}.
\section*{Data and code availability}
The accompanying repository contains claim registries, machine-readable witnesses, independent finite regressions, and a browser-native Pyodide demonstrator. The paper distinguishes finite computational checks from the symbolic arguments proving the all-size statements.
\section*{Declarations}
The authors declare no competing interests. Author identities, affiliations, funding, and acknowledgments are intentionally left blank for journal submission preparation.
\nocite{*}
\bibliographystyle{plain}
\bibliography{references}
\end{document}
